\begin{lemma}[Proved: full Kummer \'etale finite generation on a proper log scheme]
\label{D:proper-log-finite-generation}
Let $k$ be algebraically closed and let $X^{\log}$ be a nonempty
connected fine saturated log scheme whose underlying scheme $X$
is proper over $k$.  Then its full Kummer \'etale fundamental
group is topologically finitely generated.
\end{lemma}
\begin{proof}
We first check that every finite Kummer \'etale cover is tame
relative to $k$ in the precise valuative sense of
\cite[Definition 3.2.4]{AHLS2026}.  Let $(K,V)$ be a strictly
henselian valued field over $k$, and let
$\operatorname{Spec}K\to X$ be a test point.  Properness extends
this map to $\operatorname{Spec}V\to X$.  Equip
$T=\operatorname{Spec}V$ with the pulled-back fine saturated
log structure.  It has a global fine saturated chart: an
\'etale chart near the closed point pulls back along a section
over the strictly henselian local scheme, and an open
neighborhood of its closed point is the entire scheme.

By \cite[Proposition 3.1.8]{AHLS2026}, the full Kummer \'etale
fundamental group of $T$ is a free module of finite rank over
$\widehat{\mathbb Z}'$, with the residue characteristic omitted.
In particular every finite Kummer \'etale cover of $T$ is a
prime-to-residue-characteristic object of its Galois category.
Pullback to $\operatorname{Spec}K$ preserves this property.
This is exactly the condition in the cited definition, for
every test valued field.  Thus
\[
 \operatorname{FEt}^{\mathrm{ket}}(X^{\log})
   =\operatorname{FEt}^{\mathrm{t,ket}}(X^{\log}/k).
\]
Every fine saturated log scheme is of type $V_d$ in the
terminology of that source.  Its Theorem 3.5.1 therefore makes
the group on the right finitely generated, and proves the
assertion for the full Kummer \'etale group.
\end{proof}

\begin{proposition}[Proved: fixed-boundary finite generation for proper log regular pairs]
\label{D:log-regular-boundary}
Let $X^{\log}$ be a connected fine saturated log regular scheme
whose underlying scheme $X$ is proper over an algebraically
closed field.  Let $U$ be its interior, where the log structure
is trivial, and put $Y=X\setminus U$.  Then the group $Q(X,Y)$,
defined using tameness only at the original codimension-one
points of $Y$, is topologically finitely generated.
Neither projectivity nor regularity of the underlying $X$ is
assumed.
\end{proposition}
\begin{proof}
The underlying scheme of a log regular scheme is normal, and
its interior is dense.  Hoshi's exact comparison
\cite[Proposition A.10, pp. 57--58]{Hoshi2006B} identifies the
finite Kummer \'etale covers of $X^{\log}$ with the finite
\'etale covers of $U$ whose normalizations are tame at every
$y\in Y$ with $\dim\mathcal O_{X,y}=1$.  This is precisely the
cover category defining $Q(X,Y)$; there is no additional
exceptional-valuation condition in that comparison.
Lemma~\ref{D:proper-log-finite-generation} now proves the claim.
\end{proof}

\begin{remark}
This is a recognition criterion for the fixed compactification.
For a singular pair without a log regular structure, the
comparison with its original codimension-one tame cover category
is not supplied by the quoted theorem.  Thus this result does
not close the general fixed-boundary gate.  In particular,
placing a logarithmic structure on an arbitrary singular
compactification does not establish the required log regularity.
\end{remark}
